% This file was assembled from main.tex and manuscript inputs.
% Generated on 2026-05-30.

\documentclass{amsart}

%\usepackage[a4paper,margin=30mm]{geometry}
\usepackage{amsmath,amssymb}
\usepackage[colorlinks=true,citecolor=blue,urlcolor=blue]{hyperref}

\theoremstyle{plain}
\newtheorem{theorem}{Theorem}[section]
\newtheorem{lemma}[theorem]{Lemma}
\newtheorem{corollary}[theorem]{Corollary}

\theoremstyle{definition}
\newtheorem{definition}[theorem]{Definition}

\newcommand{\R}{\mathbb{R}}
\newcommand{\Pc}{\mathcal{P}}
\newcommand{\Lc}{\mathcal{L}}
\newcommand{\ZFC}{\mathsf{ZFC}}
\newcommand{\FMEA}{\mathsf{FMEA}}
\newcommand{\CH}{\mathsf{CH}}
\newcommand{\Con}{\operatorname{Con}}
\newcommand{\add}{\operatorname{add}}

\title{Relative Independence of Erd\H{o}s Problem \#501}
\author{Sungchul Lee}
\date{\today}

\begin{document}

\begin{abstract}
We prove a strengthening of the positive assertion in Erd\H{o}s problem \#501 assuming that Lebesgue measure extends to a measure on all subsets of $\R$.
Together with Hechler's counterexample under the continuum hypothesis, this gives a relative independence result for Erd\H{o}s problem \#501, relative to the consistency of this extension hypothesis, and hence relative to the consistency of a measurable cardinal.
\end{abstract}

\maketitle

\section{Introduction}

Let $m$ denote Lebesgue measure on $\R$, and let $m^*$ denote Lebesgue outer measure.
Erd\H{o}s problem \#501 \cite{ErdosProblems501} asks whether every family $(A_x)_{x\in\R}$ of bounded subsets of $\R$ with $m^*(A_x)<1$ admits an infinite independent set, that is, an infinite set $X\subseteq\R$ such that $x\notin A_y$ whenever $x,y\in X$ are distinct.
We write $P$ for the affirmative answer to this question.

Erd\H{o}s and Hajnal \cite{ErdosHajnal1960} proved the existence of arbitrarily large finite independent sets.
Hechler \cite{Hechler1972} proved $\neg P$ under the assumption $\CH$.
Newelski, Pawlikowski, and Seredy\'nski \cite{NPS1987} proved that if all sets $A_y$ are closed and have measure $<1$, then there is an infinite independent set.
However, it remained open whether $\neg P$ can be proved without assuming $\CH$.

Let $\FMEA$ denote the Full Measure Extension Axiom, the assertion that Lebesgue measure has an extension to a measure on all subsets of $\R$.
Our main result is the following.

\begin{theorem}\label{thm:main}
Assume $\ZFC+\FMEA$.
Every family $(A_y)_{y\in\R}$ of subsets of $\R$ with Lebesgue outer measure $<1$ admits an infinite independent set.
\end{theorem}

Note that Theorem \ref{thm:main} does not assume boundedness.
Hence Theorem \ref{thm:main} implies $P$ under $\FMEA$.
$\FMEA$ is equiconsistent with the existence of a measurable cardinal; see \cite[1D(e), 2E]{FremlinRVM}.
Combining Theorem \ref{thm:main} with Hechler's theorem gives the following consequence.

\begin{corollary}\label{cor:relative-independence}
If $\Con(\ZFC+\FMEA)$ holds, then $P$ is independent of $\ZFC$.
In particular, if $\ZFC$ with a measurable cardinal is consistent, then $P$ is independent of $\ZFC$.
\end{corollary}

The key ingredient of the proof is Theorem \ref{thm:kunen}, which gives a Fubini-type inequality for arbitrary subsets of $\R^2$.
This makes it possible to handle arbitrary sets $A_y$, without assuming, for instance, that they are closed.

\subsection*{Notation}

For a set $X$, we write $\Pc(X)$ to denote its power set.
$\R^{<\omega}$ denotes the set of all finite sequences of real numbers; a sequence $s\in\R^{<\omega}$ of length $n$ is a function from $\{0,\ldots,n-1\}$ to $\R$.

\section{Proof of the Main Theorem}\label{sec:main-theorem-proof}

This section proves Theorem \ref{thm:main} from Lemma \ref{lem:key-selection}.
We first state Lemma \ref{lem:key-selection}, postponing its proof to a later section, and then show how Theorem \ref{thm:main} follows from it.

If $\nu:\Pc(\R)\to[0,\infty]$ is a measure extending Lebesgue measure, then every set $S\subseteq\R$ satisfies
\begin{equation}\label{eq:outer-measure-domination}
\nu(S)\le m^*(S).
\end{equation}
Indeed, this follows by covering $S$ by countably many open intervals and taking the infimum over all such covers.
In particular, if $m^*(S)<1$, then $\nu(S)<1$.

For a family $(A_y)_{y\in\R}$ of subsets of $\R$, write $B_x=\{y\in\R:x\in A_y\}$ for $x\in\R$.

\begin{lemma}\label{lem:key-selection}
Let $\nu:\Pc(\R)\to[0,\infty]$ be a measure extending Lebesgue measure.
Let $(A_y)_{y\in\R}$ be a family of subsets of $\R$ such that $m^*(A_y)<1$ for every $y\in\R$.
If $C\subseteq\R$ and $\nu(C)=\infty$, then there is an $x\in C$ such that $\nu(C\setminus B_x)=\infty$.
\end{lemma}

\begin{proof}[Proof of Theorem \ref{thm:main}]
Assume $\FMEA$, and fix a measure $\nu:\Pc(\R)\to[0,\infty]$ extending Lebesgue measure.
Let $(A_y)_{y\in\R}$ satisfy $m^*(A_y)<1$ for every $y\in\R$.
Fix a well-ordering $\preceq$ of $\R$ and a point $r_0\in\R$.

For a finite sequence $s\in\R^{<\omega}$, define
\begin{equation}\label{eq:definition-of-Cs}
C_s
=
\R\setminus\bigcup_{i<|s|}\left(A_{s(i)}\cup B_{s(i)}\cup\{s(i)\}\right),
\end{equation}
where $|s|$ denotes the length of $s$.
Let
\begin{equation}\label{eq:definition-of-Qs}
Q_s
=
\{a\in C_s:\nu(C_s\setminus B_a)=\infty\}.
\end{equation}
Define $G:\R^{<\omega}\to\R$ by
\[
G(s)=
\begin{cases}
\text{the $\preceq$-least element of $Q_s$}, & Q_s\ne\varnothing,\\
r_0, & Q_s=\varnothing.
\end{cases}
\]
We write $f\upharpoonright n$ for the restriction of $f$ to $\{0,\ldots,n-1\}$.
By recursion on $\omega$, there exists $f:\omega\to\R$ such that
\[
f(n)=G(f\upharpoonright n)
\qquad(n<\omega).
\]

We prove by induction that
\[
\nu(C_{f\upharpoonright n})=\infty
\qquad(n<\omega).
\]
For $n=0$, $f\upharpoonright 0$ is the empty sequence, so $C_{f\upharpoonright 0}=\R$ and $\nu(C_{f\upharpoonright 0})=\infty$.
Suppose $\nu(C_{f\upharpoonright n})=\infty$.
Lemma \ref{lem:key-selection} applied to $C_{f\upharpoonright n}$ and \eqref{eq:definition-of-Qs} show that $Q_{f\upharpoonright n}$ is nonempty.
Hence $f(n)=G(f\upharpoonright n)$ belongs to $C_{f\upharpoonright n}$ and satisfies
\[
\nu(C_{f\upharpoonright n}\setminus B_{f(n)})=\infty.
\]
Moreover, by \eqref{eq:definition-of-Cs},
\[
C_{f\upharpoonright(n+1)}
=
C_{f\upharpoonright n}\setminus\left(A_{f(n)}\cup B_{f(n)}\cup\{f(n)\}\right)
=
(C_{f\upharpoonright n}\setminus B_{f(n)})\setminus\left(A_{f(n)}\cup\{f(n)\}\right).
\]
By \eqref{eq:outer-measure-domination},
\[
\nu(A_{f(n)})\le m^*(A_{f(n)})<1,
\]
and $\nu(\{f(n)\})=0$ because $\nu$ extends Lebesgue measure.
Thus $\nu(A_{f(n)}\cup\{f(n)\})<\infty$.
Since $\nu(C_{f\upharpoonright n}\setminus B_{f(n)})=\infty$ and $\nu(A_{f(n)}\cup\{f(n)\})<\infty$, we have $\nu(C_{f\upharpoonright(n+1)})=\infty$.
This completes the induction.

Set $X=\{f(n):n<\omega\}$.
Let $i<j$.
By \eqref{eq:definition-of-Cs},
\[
f(j)\in C_{f\upharpoonright j}\subseteq C_{f\upharpoonright(i+1)}
=C_{f\upharpoonright i}\setminus\left(A_{f(i)}\cup B_{f(i)}\cup\{f(i)\}\right),
\]
we have $f(j)\ne f(i)$, $f(j)\notin A_{f(i)}$, and $f(j)\notin B_{f(i)}$.
The last relation is equivalent to $f(i)\notin A_{f(j)}$.
Thus $X$ is infinite, and $x\notin A_y$ for all distinct $x,y\in X$.
\end{proof}

\section{Proof of Lemma \ref{lem:key-selection}}\label{sec:key-selection}

We use the following theorem of Kunen \cite[543C]{FremlinVol5Ch54}.
For a topological space $X$, $w(X)$ is the least cardinality of a base for the topology of $X$.
For a measure $\nu$, define $\add(\nu)$ as in \cite[511G(a)]{FremlinVol5Ch51}.
\begin{theorem}\label{thm:kunen}
Suppose that $(Y,\Pc(Y),\nu)$ is a $\sigma$-finite measure space and that $(X,\mathcal T,\Sigma,\mu)$ is a $\sigma$-finite quasi-Radon measure space\footnote{See \cite[411H(a)]{FremlinVol4Ch41} for the definition.} with $w(X)<\add(\nu)$.
Then every function $F:X\times Y\to[0,\infty]$ satisfies
\begin{equation}\label{eq:fremlin-kunen}
\overline{\int_X}\int_Y F(x,y)\,d\nu(y)\,d\mu(x)
\le
\int_Y \overline{\int_X} F(x,y)\,d\mu(x)\,d\nu(y).
\end{equation}
\end{theorem}

For the measure $\nu$ in Lemma \ref{lem:key-selection}, $(\R,\Pc(\R),\nu)$ is $\sigma$-finite and $(\R,\Lc,m)$ is a $\sigma$-finite quasi-Radon measure space, where $\Lc$ denotes the $\sigma$-algebra of Lebesgue measurable subsets of $\R$.
Moreover, $\R$ has a countable base, while countable additivity of $\nu$ gives
\[
w(\R)=\omega<\add(\nu).
\]
Applying Theorem \ref{thm:kunen} to the indicator of any $H\subset\R^2$, we get
\begin{equation}\label{eq:section-inequality}
\overline{\int_\R}\nu(H_x)\,dm(x)
\le
\int_\R m^*(H^y)\,d\nu(y),
\end{equation}
where
\[
H_x=\{y\in\R:(x,y)\in H\},
\qquad
H^y=\{x\in\R:(x,y)\in H\}.
\]
Here $\overline{\int_\R} g\,dm$ denotes the Lebesgue upper integral; that is, for $g:\R\to[0,\infty]$,
\[
\overline{\int_\R} g\,dm
=
\inf\left\{\int_\R h\,dm:g\le h,\ h\text{ is Lebesgue measurable}\right\}.
\]

We now prove Lemma \ref{lem:key-selection}.
Let $C\subseteq\R$ satisfy $\nu(C)=\infty$.
Suppose, toward a contradiction, that
\begin{equation}\label{eq:contrary}
\nu(C\setminus B_x)<\infty,
\qquad
\forall x\in C.
\end{equation}

For $N\ge1$, put $C_N=C\cap[-N,N]$.
Since $C=\bigcup_{N\ge1}C_N$ and the sequence $(C_N)_{N\ge1}$ is increasing, continuity from below gives
\begin{equation}\label{eq:CN-to-infinity}
\nu(C_N)\longrightarrow\infty.
\end{equation}
Choose $N$ such that $D:=C\cap[-N,N]$ satisfies $\nu(D)>1$.
For $k\ge1$, define
\begin{equation}\label{eq:definition-of-Dk}
D_k=\{x\in D:\nu(C\setminus B_x)\le k\}.
\end{equation}
By \eqref{eq:contrary}, $D=\bigcup_{k\ge1}D_k$, and the sequence $(D_k)_{k\ge1}$ is increasing.
Again by continuity from below, there is $k\ge1$ such that $\nu(D_k)>1$.
By \eqref{eq:outer-measure-domination},
\begin{equation}\label{eq:Dk-large-outer}
m^*(D_k)>1.
\end{equation}
Also $D_k\subseteq[-N,N]$, so
\begin{equation}\label{eq:Dk-finite-outer}
m^*(D_k)<\infty.
\end{equation}

For $M>N$,
\[
\nu(C_M)\le\nu([-M,M])=2M<\infty.
\]
By \eqref{eq:CN-to-infinity}, let $M>N$ be large enough that $\nu(C_M)>k$.

If $x\in D_k$, then $\nu(C\setminus B_x)\le k$ by \eqref{eq:definition-of-Dk}.
Since $C_M$ has finite $\nu$-measure and $C_M\setminus B_x\subseteq C\setminus B_x$, we have
\begin{equation}\label{eq:vertical-lower-section}
\nu(B_x\cap C_M)
=
\nu(C_M)-\nu(C_M\setminus B_x)
\ge
\nu(C_M)-k,
\qquad
\forall x\in D_k.
\end{equation}

Let
\[
H=\{(x,y)\in D_k\times C_M:x\in A_y\}\subseteq\R^2.
\]
For $x\in\R$,
\[
H_x=
\begin{cases}
B_x\cap C_M, & x\in D_k,\\
\varnothing, & x\notin D_k.
\end{cases}
\]
Thus \eqref{eq:vertical-lower-section} implies
\[
\nu(H_x)\ge(\nu(C_M)-k)1_{D_k}(x)
\qquad(x\in\R).
\]
By monotonicity of upper integral and the identity
\[
\overline{\int_\R} a1_S\,dm=a\,m^*(S)
\qquad(a\ge0,
\ S\subseteq\R),
\]
we get
\begin{equation}\label{eq:left-lower}
\overline{\int_\R}\nu(H_x)\,dm(x)
\ge
(\nu(C_M)-k)m^*(D_k).
\end{equation}

For $y\in\R$,
\[
H^y=
\begin{cases}
A_y\cap D_k, & y\in C_M,\\
\varnothing, & y\notin C_M.
\end{cases}
\]
Since $m^*(A_y)<1$ for every $y$,
\[
m^*(H^y)
\le
1_{C_M}(y)
\qquad(y\in\R).
\]
The function $y\mapsto m^*(H^y)$ is $\nu$-measurable because $\nu$ is defined on all subsets of $\R$.
Hence
\begin{equation}\label{eq:right-upper}
\int_\R m^*(H^y)\,d\nu(y)
\le
\int_\R 1_{C_M}(y)\,d\nu(y)
=
\nu(C_M).
\end{equation}

Applying \eqref{eq:section-inequality} to this $H$ and using \eqref{eq:left-lower} and \eqref{eq:right-upper}, we obtain
\[
(\nu(C_M)-k)m^*(D_k)
\le
\nu(C_M).
\]
Dividing by $\nu(C_M)>0$ gives
\[
\left(1-\frac{k}{\nu(C_M)}\right)m^*(D_k)
\le
1.
\]
Letting $M\to\infty$ through integers greater than $N$ and using \eqref{eq:CN-to-infinity}, together with \eqref{eq:Dk-finite-outer}, yields $m^*(D_k)\le1$.
This contradicts \eqref{eq:Dk-large-outer}.
Therefore \eqref{eq:contrary} is false, and some $x\in C$ satisfies $\nu(C\setminus B_x)=\infty$.
This proves Lemma \ref{lem:key-selection}.

\appendix

\section{Counterexample Under CH}\label{sec:ch-counterexample}

For completeness, we give a simple proof that $P$ fails under $\CH$.
Assume $\CH$, and fix an enumeration $\R=\{r_\alpha:\alpha<\omega_1\}$.
Let $\preceq$ be the induced well-ordering of $\R$, so that $r_\alpha\preceq r_\beta$ if and only if $\alpha\le\beta$, and let $\prec$ be the corresponding strict order.
For every $\beta<\omega_1$, the set $\{r_\alpha:\alpha<\beta\}$ is countable.
For $y=r_\beta$, define
\[
A_y
=
\{r_\alpha:\alpha<\beta,\ |r_\alpha|\le |y|+1\}.
\]
Then $A_y$ is countable, and hence $m^*(A_y)=0$.
Also $A_y\subseteq[-(|y|+1),|y|+1]$, so $A_y$ is bounded.

We claim that the family $(A_y)_{y\in\R}$ has no infinite independent set.
Suppose otherwise, and let $X\subseteq\R$ be infinite and independent.
Choose a strictly increasing sequence
\[
x_0\prec x_1\prec x_2\prec\cdots
\]
from $X$.
If $i<j$, then independence gives $x_i\notin A_{x_j}$.
Since $x_i\prec x_j$, the definition of $A_{x_j}$ implies
\[
|x_i|>|x_j|+1.
\]
In particular, $|x_n|>|x_{n+1}|+1$ for every $n<\omega$.
Iterating, we get $|x_n|<|x_0|-n$, which is impossible for $n>|x_0|$.
Therefore no infinite independent set exists, and $P$ fails under $\CH$.

\section*{Use of AI}

During the preparation of this work, the author used OpenAI's GPT-5.5 Pro to explore possible proof strategies and to obtain the main methodological ideas behind the argument.
After using this tool, the author independently verified the mathematical details, refined the exposition, and takes full responsibility for the final manuscript.

\begin{thebibliography}{99}

\bibitem{ErdosProblems501}
T. F. Bloom,
Erd\H{o}s Problem \#501,
\url{https://www.erdosproblems.com/501},
accessed 29 May 2026.

\bibitem{ErdosHajnal1960}
P. Erd\H{o}s and A. Hajnal,
Some remarks on set theory. VIII,
\emph{Michigan Math. J.} \textbf{7} (1960), 187--191.

\bibitem{FremlinVol5Ch54}
D. H. Fremlin,
\emph{Measure Theory}, Vol. 5, Chapter 54, Real-valued-measurable cardinals,
version of 23 October 2014 / 17 September 2016,
\url{https://www1.essex.ac.uk/maths/people/fremlin/chap54.pdf}.

\bibitem{FremlinVol5Ch51}
D. H. Fremlin,
\emph{Measure Theory}, Vol. 5, Chapter 51, Cardinal functions,
version of 9 January 2015,
\url{https://www1.essex.ac.uk/maths/people/fremlin/chap51.pdf}.

\bibitem{FremlinVol4Ch41}
D. H. Fremlin,
\emph{Measure Theory}, Vol. 4, Chapter 41, Topologies and measures I,
version of 14 December 2006,
\url{https://www1.essex.ac.uk/maths/people/fremlin/chap41.pdf}.

\bibitem{FremlinRVM}
D. H. Fremlin,
\emph{Real-valued-measurable cardinals},
version of 19 September 2009,
\url{https://www1.essex.ac.uk/maths/people/fremlin/rvmc.pdf}.

\bibitem{Hechler1972}
S. H. Hechler,
On two problems in combinatorial set theory,
\emph{Bull. Acad. Polon. Sci. S\'er. Sci. Math. Astronom. Phys.} \textbf{20} (1972), 429--431.

\bibitem{NPS1987}
L. Newelski, J. Pawlikowski, and W. Seredy\'nski,
Infinite free set for small measure set mappings,
\emph{Proc. Amer. Math. Soc.} \textbf{100} (1987), 335--339.

\end{thebibliography}


\end{document}
